\section{Recorded numerical support}\label{sec:support}

The finite analytic framework has been applied to one actual regular
theta fiber, without an RH assumption. The recorded outward-ball
calculation \cite{ThetaRecord}, using Arb arithmetic \cite{Arb}, treats
$t=0$ and $|z-35|\le8$. A complete boundary winding count certifies
exactly two interior roots, each separately enclosed as simple and real:
\[
 \alpha\simeq28.26945028346938758091450396712,\qquad
 \beta\simeq42.04407927754310998525695918780.
\]
The certificate uses the exact normalization \eqref{eq:xi}, with complete
analytic Taylor and divided-difference tails, rather than a finite source
head treated as the full field. Its native boundary interval is
\begin{equation}\label{eq:numBoundary}
 [28770927.163097,\ 33752291.350598],
\end{equation}
and its native direct-coordinate interval is
\begin{equation}\label{eq:numDirect}
 [28728270.232214,\ 34655897.029090].
\end{equation}
The endpoints are outward rational bounds as printed. Both intervals
meet the recorded common relative-diameter target, and the first lies
strictly inside the second. The result supports a finite acquisition
on a genuine regular fiber. It compares the two specified native
implementations; a direct method permitted to reuse all the same
identities and observations inherits their intersection. Neither the
interval containment nor the recorded timing difference establishes an
intrinsic advantage for the energy method. Detailed arithmetic,
accounting, and replay instructions remain in \cite{ThetaRecord}.
No theorem above uses these numerical intervals.

\section{What has been achieved and what remains open}\label{sec:open}

\subsection{The present analytic conclusion}
The project now has a complete finite-domain identity, a full quadratic
calibration, explicit bounds for the energy and both boundary payments
after quadratic subtraction, and a source-controlled ordinary-event
chart. The source tail estimates in Appendix~\ref{sec:tails} make the
constants effective without replacing the complete canonical source.
The sharper upper bound in Section~\ref{sec:heat} improves a specified
coarse bound on an explicit event-dependent interval. These statements
hold with all spatial boundaries, both sheets, the mixed term, and the
relative period retained.

The same analysis explains its current limitation. At an ordinary event,
\[
 \E_\tau=\frac{2\pi}{|A|}\log(1/\tau)+O(1),
\]
and the derived upper bounds have this permitted growth. Their constants
are finite on each fixed chart but are not a theta-specific obstruction
to that chart. The polynomial heat family $z^2-2\tau$ satisfies the local
mechanism exactly. Consequently ordinary holomorphicity, the heat
equation, real symmetry, and a controlled nonvanishing unit cannot alone
supply the desired exclusion.

\subsection{Endpoint control and coverage}
For a collision exclusion through reciprocal energy, one would need a
new estimate for the unchanged complete source, for example
\begin{equation}\label{eq:missingbound}
 \sup_{0<\tau\le\tau_0}\E(H(T+\tau,\cdot);D)<\infty
\end{equation}
on a fixed disk containing the collision, with hypotheses not already
assuming the conclusion. The logarithmic lower bound would then give
a contradiction. No estimate of the form \eqref{eq:missingbound} is
proved here. Alternatively, a source inequality that forces a sufficiently
strong upper control of the relative period, or of the full exterior
payment after the common constant is fixed, could address the same gap.
It must be checked against the complete quadratic calculation.

Even such an ordinary-event estimate would need a coverage argument.
The present theorem assumes two real simple roots in a fixed disk and
starts from a hypothetical real nondegenerate double event. It does not
treat every height, higher multiplicities, arbitrary nonreal pairs,
or a moving spatial boundary. An individual simple nonreal zero still
has finite reciprocal energy by Proposition~\ref{prop:multi}; regular
fiber finiteness by itself does not locate zeros on the real axis.
Any implication toward RH or toward RH plus simplicity therefore needs
an additional global bridge, with its reality regime, endpoint passage,
and uniform constants stated separately. The usual canonical heat-flow
context is described in \cite{RT}; it does not provide the missing
energy estimate merely by a change of language.

\subsection{Controls on a proposed source-specific advance}
A future source-specific inequality should state exactly which canonical
properties it uses and test their separate roles. A positive finite
weighted-square or finite-mode control may preserve useful positivity
while failing the complete normalized reflection identity and the
unchanged canonical coefficients. A reflected primitive control may
preserve evenness or a transform identity while losing the canonical
square-index sequence. An Euler-product modification may change the
arithmetic source itself. Such controls answer different questions;
none should be said to preserve the entire canonical class merely
because one identity survives.

The relevant next analytic task is thus to isolate a complete-source
input that controls degeneration beyond the generic chart bounds.
The finite identities and paid remainders here specify where that input
must enter. A compact-surface or Jensen-based reserve approach would
likewise require its own quantitative passage from finite bordered
surfaces; the compact bilinear relation alone does not remove the
physical boundary or supply the needed uniformity.
